% 2. Une discussion sur les changements à apporter pour passer de la version prototype
% à la version commerciale, tel que mentionné dans le mandat. Incluant une discussion
% sur le microcontrolleur, les interfaces des capteurs et une discussion sur les réseaux
% sans �l.

\subsection{Choix du microcontrolleur}

\todo{faire un choix et l'expliquer}
\todo{justifier avec le contexte de la problématique
(c'est ok d'inventer quand la problématique dit rien)}
\todo{faire des liens avec les autres domaines: controlleur+capteur+sensor
(obligatoire: le prof pas du cas-par-cas)}


\subsection{Choix d'interfaces capteurs}

\todo{faire un choix et l'expliquer}
\todo{justifier avec le contexte de la problématique
(c'est ok d'inventer quand la problématique dit rien)}
\todo{faire des liens avec les autres domaines: controlleur+capteur+sensor
(obligatoire: le prof pas du cas-par-cas)}


\subsection{Choix de communications réseaux}

\todo{faire un choix et l'expliquer}
\todo{justifier avec le contexte de la problématique
(c'est ok d'inventer quand la problématique dit rien)}
\todo{faire des liens avec les autres domaines: controlleur+capteur+sensor
(obligatoire: le prof pas du cas-par-cas)}
